\documentclass[journal,onecolumn]{IEEEtran}

% ── Encabezado ──────────────────────────────────────────────
\usepackage{fancyhdr}
\pagestyle{fancy}
\fancyhf{}
\renewcommand{\headrulewidth}{0.4pt}
\fancyhead[L]{\small Conclusiones}
\fancyfoot[C]{\thepage}

% ============================================================
% IDIOMA Y CODIFICACIÓN  (mover ANTES de csquotes)
% ============================================================
\usepackage[utf8]{inputenc}
\usepackage[T1]{fontenc}
\usepackage[spanish,es-tabla]{babel}

% ============================================================
% CITAS Y TEXTO
% ============================================================
\usepackage{csquotes}

% ============================================================
% ESPACIADO Y LISTAS
% ============================================================
\usepackage{parskip}
\setlength{\parskip}{4pt}
\usepackage{enumitem}

% ============================================================
% IDIOMA Y CODIFICACIÓN
% ============================================================
\usepackage[spanish,es-tabla]{babel}
\usepackage[utf8]{inputenc}
\usepackage[T1]{fontenc}

% ============================================================
% PAPEL Y GEOMETRÍA
% ============================================================
\usepackage{geometry}
\geometry{a4paper, top=2.5cm, bottom=2.5cm, left=2.5cm, right=2.5cm}

% ============================================================
% IMÁGENES
% ============================================================
\usepackage{graphicx}
\graphicspath{{Figuras/}}

% ============================================================
% TABLAS
% ============================================================
\usepackage{booktabs}
\usepackage{tabularx}
\usepackage{longtable}

% ============================================================
% MATEMÁTICAS
% ============================================================
\usepackage{amsmath}
\usepackage{amssymb}

% ============================================================
% CÓDIGO FUENTE
% ============================================================
\usepackage{listings}

% ============================================================
% FORMATO ACADÉMICO
% ============================================================
\usepackage{microtype}
\usepackage{float}

% ============================================================
% CAPTIONS (TÍTULOS DE FIGURAS/TABLAS)
% ============================================================
\usepackage{caption}
\captionsetup{font=small, labelfont=bf}

% ============================================================
% REFERENCIAS / HIPERVÍNCULOS
% ============================================================
\usepackage[hidelinks]{hyperref}

% ============================================================
% NUMERACIÓN Y NOMBRES EN ESPAÑOL
% ============================================================
\renewcommand{\thetable}{\arabic{table}}
\def\tablename{Tabla}

\renewcommand{\thesection}{\arabic{section}}
\renewcommand{\thesubsection}{\thesection.\arabic{subsection}}
\renewcommand{\thesubsubsection}{\thesubsection.\arabic{subsubsection}}

% ============================================================
% ESTILO DE TÍTULOS DE SECCIÓN
% ============================================================
\usepackage{titlesec}
\titleformat{\section}{\normalfont\Large\bfseries}{\thesection}{1em}{}
\titleformat{\subsection}{\normalfont\large\bfseries}{\thesubsection}{1em}{}
\titleformat{\subsubsection}{\normalfont\normalsize\bfseries}{\thesubsubsection}{1em}{}

% ============================================================
% CADA SECCIÓN PRINCIPAL (\section) INICIA EN PÁGINA NUEVA
% ============================================================
\usepackage{etoolbox}
\pretocmd{\section}{\clearpage}{}{}

% ============================================================
% BIBLIOGRAFÍA (BIBLATEX + BIBER)
% ============================================================
\usepackage[
backend=biber,
style=ieee,
sorting=none
]{biblatex}
\addbibresource{referencias_C1.bib}

% ============================================================
% DOCUMENTO
% ============================================================
\begin{document}

\subsection*{Conclusión 1 --- Unidad 1: Fundamentos y elicitación}
\label{sec:conclusion1}

La Unidad 1 del curso combina dos temas que el equipo trabajó en prácticas
distintas: en la PE1 se construyeron los fundamentos y el contexto del
sistema mediante análisis documental, sin elicitación de campo todavía; en
la PE2 se ejecutó la elicitación real, con entrevistas y un cuestionario
digital. Por eso la evidencia de esta conclusión proviene de la PE2, que es
la práctica donde efectivamente se puede responder qué técnica reveló
requisitos nuevos y qué actor quedó sin representación directa en el
proceso.

La técnica que reveló requisitos inéditos no fue el cuestionario (EV-07),
que solo cuantificó preferencias ya conocidas, sino la combinación de
entrevista con sesión de \emph{walkthrough}. Las dos últimas evidencias
recolectadas, EV-23 y EV-24, produjeron directamente RF-27 (control de
cruzas repetitivas, por riesgo de maltrato animal) y RF-26 (segunda opinión
veterinaria), ninguno anticipado por las quince entrevistas previas ni por
el cuestionario \cite{pohl2025}. Además, el Refugio de animales
clasificado como actor clave en la matriz de poder e interés de la PE1
no fue entrevistado, encuestado ni observado en ninguna de las 24
evidencias de campo; el equipo compensó esto de forma indirecta, trazando
sus requisitos contra evidencia de propietarios y observación, sin fuente
propia del refugio \cite{swebok2024}.

Esto muestra que la calidad de la elicitación no depende de cuántas
técnicas se apliquen, sino de qué tan bien cada una llega a la fuente
correcta: el cuestionario amplió cobertura pero no profundidad, mientras
que la conversación abierta con revisión guiada del sistema fue la única
capaz de sacar a la luz necesidades que los propios usuarios no habían
verbalizado antes de interactuar con él. Al mismo tiempo, la ausencia del
Refugio de animales deja claro que un stakeholder puede quedar
``cubierto'' en la matriz de trazabilidad sin haber sido consultado nunca
directamente, lo cual es una brecha real de validez del ERS no grave
por sí sola, pero sí una limitación que el equipo debe reconocer
explícitamente en vez de presentar la cobertura indirecta como
equivalente a una fuente primaria.



\printbibliography[heading=none]

\end{document}
